% =========================================================================
% PPCF draft (single column). Main text, rebuilt a third time on 2026-11-18 under the
% author's three rules: a sentence stays only if the claim or the reader needs it; numbers
% go into tables and figures, not into running text (the abstract keeps five); length comes
% from figures and tables, not from words.
% Claim: on the KSTAR edge CES chord selected as the pedestal top, the beam-emission value
% of the present frame carries information about the reported ion temperature of that
% frame beyond the CES history; measured three ways (recurrent network against its twin,
% five-parameter filter with and without a BES drive, the trained network given
% one-frame-old inputs) and sized against a floor set by the scatter of the reported value.
% The paper's estimator is the two-stage estimator (B.64, r2-r4: the network as first
% stage plus a boosted second stage on its error); the mechanism analyses read the first
% stage. Everything else (four-split ledger, average of three initialisations, measurement
% regime and MAST radial test, the first readouts with whole-discharge standardization,
% thinned replay, edge bursts, MAST pedestal top, rotation, the 20-parameter reading, the
% earlier second stages on single splits, the wall-field note of B.63, registration
% history) is in supp_ppcf.tex.
% Numbers of the tables and of the macros are generated:
%   py docs/paper/ppcf/make_tables.py  (b39_primary, b54_pooled_rows, b55_sample_timing,
%   b56_kfx, b57_yardsticks, b60_relative_error, b61_*, b53_components, b64_two_stage json)
% Figures: py docs/paper/ppcf/fig_paper.py (fig_geometry, fig_schematic, fig_skill,
%   fig_where, fig_blips, fig_use) and fig_trace_causal.py (fig_trace_causal)
% Check: py docs/paper/ppcf/check_numbers.py
% Build (from this folder): latexmk -pdf supp_ppcf, then latexmk -pdf main_ppcf (the two
% files cross-reference each other through xr)
% =========================================================================
\documentclass[11pt,a4paper]{article}

\usepackage[margin=2.5cm]{geometry}
\usepackage[T1]{fontenc}
\usepackage[utf8]{inputenc}
\usepackage{newtxtext,newtxmath}
\usepackage{amsmath,bm}
\usepackage{graphicx}
\usepackage{booktabs}
\usepackage{multirow}
\usepackage[skip=6pt]{caption}
\usepackage{microtype}
\usepackage{flafter}
\usepackage[numbers,sort&compress]{natbib}
\setlength{\bibsep}{0.5pt}
\renewcommand{\bibfont}{\footnotesize}
\usepackage{xr-hyper}
\usepackage[colorlinks=true,linkcolor=blue,citecolor=blue,urlcolor=blue]{hyperref}
\externaldocument[S-]{supp_ppcf}

\newcommand{\Ti}{\ensuremath{T_i}}
\newcommand{\Vrot}{\ensuremath{V_\mathrm{rot}}}
\newcommand{\Supp}[1]{Supplementary \S\ref{S-#1}}
\newcommand{\SuppTab}[1]{Supplementary Table~\ref{S-#1}}
\newcommand{\SuppFig}[1]{Supplementary Figure~\ref{S-#1}}
\input{b39_generated}

\hypersetup{
  pdftitle={Causal nowcasting of the charge-exchange ion temperature on an edge chord of KSTAR with co-located fast diagnostics},
  pdfauthor={Seungsang Lee},
}

\title{Causal nowcasting of the charge-exchange ion temperature on an edge chord of KSTAR
with co-located fast diagnostics}
\author{Seungsang Lee\\
  {\small Department of Nuclear Engineering, Seoul National University, Seoul, Republic of Korea}\\
  {\small\texttt{lss010330@snu.ac.kr}}}
\date{}

\begin{document}
\maketitle

% =========================================================================
\begin{abstract}
\noindent
Charge-exchange spectroscopy (CES) on KSTAR reports the temperature of C\,VI ions every
10\,ms; on the edge chord selected as the pedestal top the signal is weak and values are
missing. We ask what the beam-emission (BES) value of the present frame adds to the CES
history when the reported temperature of that frame is estimated causally. In a
cross-validation pooled over 639 discharges, with values above 3\,keV treated as missing,
a recurrent network that also reads BES, electron-cyclotron-emission imaging and Mirnov
channels has a mean-squared-error skill of \Ptwo{} over the same network trained without
these channels. A five-parameter robust filter driven by two summaries of the BES array
gains \KXone{} over the same filter without the drive, and the network loses its gain
when BES is one frame old: the information is in the BES value of the present frame. Most
of the network's correction is still present in the next, independently fitted frame.
With a second stage, a boosted regression of the network's error fitted on out-of-fold
errors of the non-test discharges, our estimator has on three of the eight re-partitions a skill of \TwoGp{} over
a causal Gaussian process on the CES history, a root-mean-square error lower by
\TwoRmseAbs\,eV, a median absolute error below that of holding the last value, more frames
within 20\% of the reported value than that process, and a gain over the same two stages
without the fast channels. For the network alone the gain is small against the
frame-to-frame scatter of the reported value, is not established in the coldest third of
the discharges by temperature, rests on the tenth of them with the largest gain and was
not established on single file-level splits. Use requires BES of the present frame and a
CES value at most two frames late; real-time control is a prospect.
\end{abstract}

\noindent\textbf{Keywords:} charge-exchange spectroscopy, ion temperature, pedestal,
real-time estimation, beam-emission spectroscopy, recurrent neural network, KSTAR

% =========================================================================
\section{Introduction}
\label{sec:intro}

Ion temperature \Ti{} is measured in a tokamak by charge-exchange spectroscopy
(CES)~\citep{Ko2010ChargeExchange}, from the Doppler width of a line that an impurity
ion emits after charge exchange with the atoms of a neutral beam. The measurement
integrates photons over an exposure and then fits a line shape, so it is slow and fails
intermittently, most often where the signal is weakest. The KSTAR chord studied here was
selected by the data provider as the pedestal top within the beam-emission coverage
($R = 2.198$\,m). Its 10\,ms frames lack a \Ti{} value more often when the beam emission
is low, so the values are not missing completely at
random~\citep{Little2002StatisticalAnalysis}. Beam-emission spectroscopy (BES) and
electron-cyclotron-emission imaging (ECEI) view the same location without gaps; neither
measures \Ti{}.

Two premises stand behind the question. \emph{What the value is.} CES reports the
temperature of the C\,VI impurity as fitted. In the core it has been cross-checked
against other measurements, on KSTAR against x-ray crystal spectroscopy, which does not
reach the edge~\citep{Lee2020CrossComparisons}. In the pedestal the impurity and
main-ion temperatures can differ: at ASDEX Upgrade the deuterium temperature exceeded the
impurity temperature at the pedestal top when the electron heating was
raised~\citep{CanoMegias2025MultiIon}, and at DIII-D the main-ion temperature fell to
about half the impurity temperature near the separatrix~\citep{Haskey2018MainIon}. \emph{Why a value per
frame.} Requirement documents ask for the ion temperature at 10 to 100\,ms and about
10\%: the 10\,ms figure for the edge is from the specifications developed in the
FIRE design study (Fusion Ignition Research Experiment, a proposed burning-plasma
device; p.~288 and table~III of~\citep{Young2008PlasmaMeasurements}), and the ITER
requirement is 100\,ms and 10\% for the edge ion
temperature~\citep{Donne2007ChapterDiagnostics} and 10 to 30\% for the
CXRS~\citep{Mazon2025Diagnostics}. Real-time charge-exchange analysis
exists at DIII-D, sent to the plasma control system and less accurate on edge
views~\citep{Chrystal2026RealTime}, and at JT-60U, used for feedback on the
ion-temperature gradient (\citep{Yoshida2009RealTime} and p.~S1049-2
of~\citep{Kobayashi2007DevelopmentRealtime}). On
KSTAR the conventional analysis waits for a beam-off frame for its background; a
two-Gaussian fit gives \Ti{} without beam modulation~\citep{Lee2020TwoGaussian}, and a
neural-network analysis of KSTAR spectra reports a decline of its accuracy at
$R \gtrsim 2.2$\,m and names real-time CES as its goal~\citep{Lee2026NeuralNetworkbased}.

An estimate for monitoring or control has to be formed when the frame is taken, without
the next value, which may also arrive late; the simplest reference is to hold the last
value. We ask whether fast diagnostics at the location of the chord lower the error of
such a causal estimate of the reported \Ti{} below what the CES history alone allows, how
large that is against the error that can be removed at all, which diagnostic carries it,
in which frames, and under which conditions of use. In short, the BES value of the
present frame carries information about the reported \Ti{} of that frame beyond the CES
history; two estimators of different kinds, a network against its twin and a
five-parameter filter, agree on this, and the network given one-frame-old BES locates it
in the present frame (\S\ref{sec:res:main}, \S\ref{sec:res:diag}); with a second stage
fitted on the network's errors, the fast-channel increment keeps an interval above zero
and the ratio of the median error to that of hold has an interval below one
(\S\ref{sec:res:main}); for the network,
the information is small against the scatter of the reported value, is not established in
the coldest third of the discharges and mostly lasts into the next frame (\S\ref{sec:res:where},
\S\ref{sec:res:use}).

\paragraph{Relation to earlier work.} Gaussian-process (GP)
regression~\citep{rasmussen2006} is used for fusion profile
fitting~\citep{chilenski2015,michoski2024} and has been applied to KSTAR
CES~\citep{Kim2024KineticProfile}. Neural networks extract \Ti{} from charge-exchange
spectra on KSTAR~\citep{Lee2026NeuralNetworkbased} and HL-3~\citep{Tian2026RealTimeCxrs};
they read a spectrum that exists, whereas the problem here is a frame for which no usable
value is in hand. Kinetic profiles are reconstructed in real
time~\citep{Shousha2023MachineLearningbased} and densified from other
diagnostics~\citep{Imrek2026TemporalSuperresolution}; at DIII-D Thomson-scattering data are
reconstructed at high time resolution from other diagnostics, without target measurements
at inference~\citep{Jalalvand2025MultimodalSuperresolution}, whereas we estimate causally
the reported value of a frame from its own history and the fast channels. Our recurrent
network takes the sparse target as the forward-filled last value, the time since it and
a mask, the simple recurrent baseline of Che et al.~\citep{Che2018RecurrentNeural}. An
extended-Kalman observer of $T_e$ at DIII-D takes \Ti{} as a multiple of
$T_e$~\citep{Morosohk2023RealTime}; we found no frame-by-frame estimator of the
charge-exchange \Ti{} of a chord. The closest precedent of our filter merges 60\,Hz
Thomson scattering with interferometry in a 1\,kHz observer on
TCV~\citep{Pastore2023ModelBased,Pastore2026ApplicationsNovel}; there the fast signal is a
measurement with a physics forward model, in our filter a drive with fitted gains, and our
update has a Huber weight~\citep{Huber1964RobustEstimation}.

The results are for one chord, one campaign and one device, with values above 3\,keV
(failed fits) treated as missing; the readouts with those values kept, a second device
(MAST) and rotation are in the supplementary material.

% =========================================================================
\section{Data and target}
\label{sec:data}
\label{sec:data:kstar}

We use 641 KSTAR discharge files of the 2022 campaign (shots 30801--32751) on a 10\,ms
grid; 639 discharges have more than one grid row and enter the splits. Each row contains
nine BES channels~\citep{Lampert2015CombinedHydrogen,KFE2026DiagnosticsDataUserGuide}, four ECEI
channels~\citep{Yun2010DevelopmentKstar}, two Mirnov-coil signals, time, and the CES values
of one outboard chord. Figure~\ref{fig:geometry} places the chord and the channels on the
major radius as far as publications allow (\Supp{sec:target}, \SuppTab{tab:plasma}). No
equilibrium is in the files, and the published locations come from different discharges
and coordinates: none gives the pedestal-top position of the 2022 discharges in metres.
A published edge analysis that includes two discharges of this set takes the edge CES
\Ti{} as a proxy for the pedestal height~\citep{Heo2025ExperimentalIdentification}.

\begin{figure}[tbp]
  \centering
  \includegraphics[width=\linewidth]{fig_geometry}
  \caption{Where the chord and the fast channels are, on the major radius $R$. Top: a
  schematic edge \Ti{} shape for orientation only (no profile of these discharges is in
  the files), the band of published pedestal-top positions of other discharges and
  campaigns, and the midplane separatrix of one 2022 discharge of this set read from a
  published figure~\citep{Lee2026SpatiotemporalStructure}. Bottom: the CES channel of the
  files and its neighbours (facility user guide), the published BES view of one discharge,
  the three BES columns of the files (inferred from the channel names, filled when
  published, hatched when inferred) and the ECEI column of the files. Positions and
  sources: \Supp{sec:target}.}
  \label{fig:geometry}
\end{figure}

\paragraph{What a row contains.} Raw BES and ECEI signals (500\,kHz; BES at 1\,MHz in one
discharge) and Mirnov signals (2\,MHz) of
twelve discharges, set against the files (\Supp{sec:target}), show that every fast-channel
value of a row is a single raw sample within one raw step of its time stamp, not an average
over the CES frame (Fig.~\ref{fig:rows}). A step of 20, 30 or 40\,ms between rows is one
to three absent frames;
each such step inside the raw deliveries coincides with a beam blip, in which the heating
beam is off and the frames exposed without beam are left out. The converse does not always
hold: in one discharge two longer interruptions removed no frame, and the files do not say
from which beam the values of those rows come. The stamp of a row lies within about 2\,ms
of the start of its 10\,ms frame (placed from the beam interruptions; \Supp{sec:target}),
so an estimate formed at the stamp uses fast-channel samples taken early in the frame it
estimates and CES values of frames that had ended. Values that repeat the previous one to
the last digit are held values, not measurements, and are removed everywhere.

\begin{figure}[tbp]
  \centering
  \includegraphics[width=\linewidth]{fig_rows}
  \caption{What a row of the 10\,ms files holds (schematic; established on the raw
  streams of twelve discharges, \Supp{sec:target}). The CES value of a row is the frame
  whose exposure starts within about 2\,ms of the stamp; the fast channels are single raw
  samples at the stamp; a frame exposed while the heating beam is off (a blip) is left
  out, so the step between rows is 20\,ms. A causal estimate for row $t$ uses the fast
  samples of rows up to $t$ and the CES values of rows before $t$.}
  \label{fig:rows}
\end{figure}

\paragraph{What the target is.} The quantity estimated is the C\,VI temperature that the
KSTAR CES analysis reports for the chord: the temperature of C$^{6+}$ ions from the
Doppler width of the charge-exchange line excited by the heating beam, the beam being
modulated for the background~\citep{Ko2010KstarCesSystem,Ko2010ChargeExchange}, which is
taken from frames in which the beam is
off~\citep{Ko2015IonTemperature,KFE2026DiagnosticsDataUserGuide}, one value per 10\,ms
frame~\citep{Kim2024KineticProfile}. Whether it equals the main-ion
temperature is not tested here. An offset that is constant in time, or a constant factor,
changes no skill of this paper, because every estimator is scored against the same
reported values; a bias that varies within a discharge would. The files carry no per-frame
fit error. Values above 3\,keV are most frequent at low beam emission (\Supp{sec:target});
taken as failed fits of the spectrum, a source of outliers in KSTAR
CES~\citep{Kim2024KineticProfile}, they are treated as missing for every estimator. The
frame-to-frame scatter of the reported value and the absence of a step at beam-off frames,
where the analysis refreshes its background, are measured in \Supp{sec:target}.
Figure~\ref{fig:trace} shows the reported value and the estimators of \S\ref{sec:methods}
in three discharges.

\begin{figure}[tbp]
  \centering
  \includegraphics[width=\linewidth]{fig_trace_causal}
  \caption{Reported \Ti{} (circles) and three causal estimates in about 0.65\,s of three test
  discharges of one split: hold (grey steps), the causal GP on the CES history (dashed)
  and the recurrent network that also reads the fast channels (red). Shaded: rows
  without a reported value; dotted: absent frames (beam blips). (a, b) Two discharges
  with many missing values and large jumps of the reported value; (c) a typical
  discharge. The network is the first stage of the estimator of \S\ref{sec:methods:est}
  (the second stage is not drawn), trained on one split with one initialisation
  (\Supp{sec:splits}).}
  \label{fig:trace}
\end{figure}

% =========================================================================
\section{Estimators and evaluation}
\label{sec:methods}

\subsection{Estimators}
\label{sec:methods:est}
Two references use the CES past of the target only: \emph{hold} (the last observed value)
and a \emph{causal GP} (Mat\'ern-3/2 kernel on the 16 most recent values, selected by
marginal likelihood at every row), the main comparison because it averages the noisy past.
Two further references use the following value and are not available in real time: PCHIP
interpolation~\citep{fritsch1980} and an offline GP with neighbours on both sides.

The \emph{recurrent network} (the network below; Fig.~\ref{fig:schematic}a;
\Supp{sec:design}) has two LSTM branches that run forward in time over every grid row. The \Ti{} branch reads the 15 fast
channels, the time step and, per target, the last observed value, its staleness and an
observed flag; a second branch estimates the rotation of the chord from the CES past only.
At grid time $t$ the network estimates the standardized \Ti{} under a masked squared-error
loss, plus a small penalty on estimates below 0\,eV. We call an estimator \emph{causal}
when it uses CES values of earlier frames, statistics of training discharges only, and
fast samples taken up to the stamp of the row. Training uses AdamW (learning rate
$10^{-3}$, weight decay $10^{-4}$, batches of 16 segments, gradient clipping at 1.0) with
early stopping on held-out training discharges, under a cap of 80 epochs. The
\emph{history-only twin} is the same network with the 15 fast channels removed from the
input of the \Ti{} branch; the skill of the network over its twin is the
\emph{fast-channel increment} of the network. Ablation networks read BES, ECEI, or BES and
ECEI beside the CES history.

\emph{The two-stage estimator}, the estimator of this paper, adds a second stage to the
network (Fig.~\ref{fig:schematic}a). At each row the second stage estimates the error that
the network makes there and adds this estimate to the network's output. It reads 15
quantities known at the row: four summaries each of the mean BES and the mean ECEI
channel over the last ten rows (the value, its departure from the ten-row mean, the
ten-row standard deviation and the change from the row before), and seven from the CES
history and the network's estimate (the time step, the age and value of the last CES
value, and the differences of the network's estimate from that value and from the causal
GP, with their absolute values); it reads no Mirnov input. It is a gradient-boosted
regression~\citep{Friedman2001GreedyFunction}: a sum of small regression trees, each
fitted to what the trees before it leave unexplained. It is fitted only on the discharges
of each fold on which the network is trained or stopped, never on test discharges, and the
network estimates it is fitted on come from networks that were not trained on those
discharges (out-of-fold estimates, as in stacking~\citep{Wolpert1992StackedGeneralization}),
so that it learns the errors the network makes on discharges it has not seen; the
protocol is in \Supp{sec:twostage}. Where the ten-row window is incomplete (the first ten
rows of a segment, about 2\% of the scored rows) the network's estimate is kept. The
second stage takes \TwoMicros\,$\mu$s per row when rows are processed together on a CPU.
One row at a time on one CPU thread, the median time is 0.49\,ms for a step of the
network, 0.64\,ms for the causal GP that the second stage reads and 1.257\,ms for a
single-row call of a second stage of this kind (\Supp{sec:twostage}). No
real-time use of BES or CES on KSTAR is published (\SuppTab{tab:plasma}), so real-time use
is a prospect. The \emph{two-stage twin} is the twin with a second stage that reads only the seven
quantities from the CES history and the twin's estimate; the skill of the two-stage
estimator over the two-stage twin is the fast-channel increment of the paper's estimator.
The analyses of which diagnostic carries the increment, where the gain is, persistence and
delay (\S\ref{sec:res:diag}--\S\ref{sec:res:use}) were made on the network alone, the first
stage, and are stated for it.

\begin{figure}[tbp]
  \centering
  \includegraphics[width=\linewidth]{fig_schematic}
  \caption{The two estimators that read the fast channels. (a) The two-stage estimator:
  the recurrent network for the \Ti{} branch (the rotation branch, which reads the CES
  history only, is not drawn; \SuppFig{fig:arch}), followed by the second stage, which
  reads summaries of the BES and ECEI values over the last ten rows, the CES history and the
  network's estimate; the twin is the same network without the fast channels.
  (b) The five-parameter filter: the state is $x = \log\Ti{}$; the prediction for a row
  is driven by the changes of two summaries of the BES array, and a CES value updates the
  state with a weight that shrinks for a large innovation $v$ of variance $S$.}
  \label{fig:schematic}
\end{figure}

\emph{A five-parameter filter} measures the increment without a network
(Fig.~\ref{fig:schematic}b; \Supp{sec:filter}). Its state is $\log\Ti{}$, a random walk
with variance $q$ per frame; an observation with variance $r$ updates it with the weight
$\min(1, k\sqrt{S}/|v|)$, so that a large departure counts less. In the prediction step the
state moves by $b_c\,\Delta c_t + b_g\,\Delta g_t$, where $c$ is the mean of the logarithms
of the nine BES channels (the common mode, in which a change of the beam appears in full),
$g$ the difference of the mean logarithms of two columns of the array (the radial
difference, in which it largely cancels), and $\Delta$ the change from the previous row,
clipped. KF0 has no drive, KFX both gains, Kc and Kg one each. Per fold, KF0 is fitted on
the training discharges first and the gains then with $q$, $r$, $k$ held. The clip and this
sequential fit were chosen on one re-partition (r2) of the pooled cross-validation. The
ratios of the filter to the floor and its persistence are read on all eight re-partitions
(r2--r9) and its comparison with the two-stage estimator on r2--r4; every other readout
of the filter is on the other seven (r3--r9).

The causal GP is the main reference because it alone meets three conditions
(Table~\ref{tab:refs}): it reads the information of our estimator without the fast
channels, under the same causal rule, and it is a widely used method whose
hyper-parameters are selected by the data at every row, not by us. The comparison is therefore not against a straw man.

\begin{table}[tbp]
  \centering
  \caption{The references and the three conditions of the main reference. Causal GP:
  Mat\'ern-3/2 kernel on the 16 most recent CES values of the segment, standardized
  within them; at every row the length scale and the noise variance are selected by the
  exact log marginal likelihood over a fixed grid of candidates.}
  \label{tab:refs}
  \footnotesize
  \begin{tabular}{lccp{0.43\linewidth}}
    \toprule
    Reference & CES values only & Causal & Who chose the settings \\
    \midrule
    causal GP & yes & yes & the marginal likelihood, per row, for the hyper-parameters; we chose the kernel, the 16 values and the candidate grid (five length scales from 5 to 80\,ms, four noise variances) \\
    hold & yes & yes & nobody: it has no settings \\
    PCHIP & yes & no & nobody: it has no settings; uses the next value \\
    offline GP & yes & no & as the causal GP, with neighbours on both sides \\
    history-only twin & yes & yes & we did: architecture and training \\
    two-stage twin & yes & yes & we did: the twin with a second stage on the CES history \\
    KF0 & yes & yes & we did: the specification; three parameters fitted on training discharges \\
    \bottomrule
  \end{tabular}
\end{table}

\subsection{Scoring, intervals and registration}
\label{sec:methods:eval}
Errors are computed in eV on rows where \Ti{} was measured and withheld from the
estimator; every estimator of a comparison is scored on the identical rows. Under the
per-discharge cap of the protocol, the scored rows of a test discharge are a seeded subset
of at most 500, drawn anew in each re-partition; in each re-partition 633 to 635 of the
639 discharges have scored rows, and four never do (\Supp{sec:rows}). The score is
the mean-squared-error skill $1 - \mathrm{MSE}_A/\mathrm{MSE}_R$~\citep{murphy1988}, given
with a 95\% interval and whether it excludes zero. The primary readout is \emph{pooled}:
shot-level 5-fold cross-validation on 639 discharges, repeated over eight re-partitions
(r2--r9; r0 and r1 belong to an earlier batch). The five test folds of a repeat are
concatenated into one skill, and a two-level bootstrap draws a repeat and then resamples
its discharges; it holds the trained networks fixed, so the fold-level interval, which
carries the training variance, is given beside it. The ablation networks of
Table~\ref{tab:carry} were retrained on every fold of r2--r4; the component networks, the
replayed networks and the earlier ablations were trained on four file-level splits of 96
test discharges (\Supp{sec:components}, \Supp{sec:splits}); every one of these networks
has one initialisation. Because a mean squared error is decided by few rows, skills under a
Huber loss and with the hundredth of the rows with the largest error of hold removed are
co-primary; because accuracy requirements are relative, we also give the share of frames
whose absolute error is within 10\% and 20\% of the reported value, and the median absolute
error.

\emph{The floor.} The floor $s^2$ is minus the covariance of successive increments of the
reported value on the \emph{triple rows}, scored rows whose previous and next frames are
scored 10\,ms apart: a lower bound of the frame-to-frame independent variance if the
increments of the underlying value are not negatively correlated (\Supp{sec:yardsticks}).
It is the mean over these rows, which a few large excursions dominate, and so larger than
the scatter of a typical row. The mean squared error over the floor is computed on the
same triple rows, where errors are smaller than on all rows, so it does not follow from
the root-mean-square error of all rows. The floor bounds what an estimate from the CES
history can remove; a present-frame signal correlated with the conditions of that frame's
fit is not bounded by it, which the persistence test of \S\ref{sec:res:use} addresses.

\emph{Registration.} For every registered batch the statements, decision rules and
expectations were committed to the author's public repository before it was run or
computed (not a third-party registry), and failed statements are reported as failed
(\Supp{app:registered}). Descriptive statistics, readings made after results were seen and
definitions fixed after an exploratory look are not registered tests and are labelled
where they occur.

% =========================================================================
\section{Results}
\label{sec:results}

\subsection{The fast-channel increment and its size}
\label{sec:res:main}

The network and its twin were trained on every fold of the pooled cross-validation and
scored with the references on the same rows (Table~\ref{tab:main}a,b, Fig.~\ref{fig:skill};
registered). The network has a lower error than every causal reference, and its margin over
its twin, the fast-channel increment of the first stage, has an interval above zero; both
margins also have intervals above zero at the fold level, under the Huber loss and with the
largest errors trimmed, so they do not rest on a few rows. In eV, the network lowers the
root-mean-square error of the causal GP by about a twentieth.

\emph{The two-stage estimator.} The second stage was fitted and read on the folds of three
re-partitions, r2--r4 (Table~\ref{tab:main}a,c, Fig.~\ref{fig:skill}a; registered, all six
statements met; \Supp{sec:twostage}). The skills of the two-stage estimator over the
network and over the causal GP have intervals above zero, the ratio of its median absolute
error to that of hold has an interval below one, and its share of frames within 20\% of the
reported value minus that of the causal GP has an interval above zero; for the network
alone neither of the last two intervals excludes its reference. Its skill over the
two-stage twin, the fast-channel increment of the paper's estimator, has an interval above
zero (beside it in Table~\ref{tab:main}a, the increment of the network on the same rows).
Added to the twin, where it reads the CES history only, the second stage also has skill
with an interval above zero, so part of its gain does not need the fast channels. Over the
offline GP, which uses the next value, the skill of the two-stage estimator has an interval
above zero and that of the network alone an interval below zero; the offline GP keeps the
lower median error and the larger share within 20\% in point value (no paired interval was
registered for these). These statements rest on three re-partitions with one initialisation per network;
the primary statements on the network, on eight re-partitions, are unchanged. The second
stage reads no information that the network does not read; it uses the same information
differently.

\begin{table}[tbp]
  \centering
  \caption{KSTAR, \Ti{}, values above 3\,keV excluded; pooled cross-validation of 639
  discharges, 95\% intervals (bold: the interval excludes zero, or one for a ratio).
  (a) Skill under squared error, under a Huber loss and with the hundredth of the rows
  with the largest error of hold removed; the network (first stage) on r2--r9, the
  two-stage estimator on r2--r4 (three re-partitions, one initialisation per network; KFX
  on these rows includes r2, on which the filter was specified), the filter rows on the
  seven re-partitions not used to choose the filter (r3--r9). (b) On all scored rows of
  r3--r9, where every estimator but the two-stage one is read again: root-mean-square error, median
  absolute error, share of frames within 20\% of the reported value (within 10\%:
  \Supp{sec:relerr}). Mean squared error over the floor (\FloorRms\,eV root mean square):
  r2--r9, triple rows only (\S\ref{sec:methods:eval}), where errors are smaller, so it does
  not follow from the RMSE column (for KF0 and KFX it includes r2, on which the filter was
  specified); its reciprocal is the share of the error that the floor
  accounts for. (c) On all scored rows of r2--r4: the two-stage estimator beside the
  network and the references; the rows differ from those of (b), so absolute errors are
  compared within a panel only. Fold-level intervals of the network over the causal GP
  and over its twin in (a):
  \PoneFold{} and \PtwoFold{}; the root-mean-square error difference of the network to the
  causal GP on r2--r9: \RmseNetMinusGp\,eV (paired differences on r2--r4:
  \SuppTab{tab:twopairs}).}
  \label{tab:main}
  \footnotesize
  \resizebox{\linewidth}{!}{\TabMain}\par\medskip
  \resizebox{\linewidth}{!}{\TabMainB}\par\medskip
  \resizebox{\linewidth}{!}{\TabMainC}
\end{table}

\begin{figure}[tbp]
  \centering
  \includegraphics[width=\linewidth]{fig_skill}
  \caption{The main result (numbers in Table~\ref{tab:main}). (a) Skill with 95\% interval
  of the two-stage estimator over each reference (r2--r4), of the network, its first stage
  (r2--r9), and of the filter family (r3--r9); filled markers: the interval excludes zero.
  The two first-stage increments, network over twin and KFX over KF0, are of the same size.
  (b) Mean squared error of each estimator over the floor $s^2$ on the triple rows of
  r2--r9, for the filters including r2, on which they were specified (bars:
  95\% interval; the two-stage estimator was not read against the floor); grey is the part
  that the floor accounts for, colour the part above it.}
  \label{fig:skill}
\end{figure}

The filter gives an increment of the same size without a network
(Fig.~\ref{fig:skill}a). KF0 is equivalent to the causal GP within $\pm0.02$ (90\%
interval): the causal GP behaves as a robust filter on the logarithm of the value, and KF0
is the reference of the filter family. The network gains over KFX about what its twin
gains over KF0, so what the network adds to the filter is a better use of the CES
history, not more information from BES.

Between two scored frames the mean squared error of every estimator lies above the floor
(Fig.~\ref{fig:skill}b; the two-stage estimator was not read against it): the floor accounts for
about three quarters of the error of the network and of the offline GP, two thirds for the
causal GP and KF0, and half for hold. Of the error above the floor, the network removes
about three tenths of what the causal GP leaves and a fifth of what its twin leaves
(\SuppTab{tab:floor}). The offline GP, which waits for the next value, has a lower error
than the network on these rows.

Within 20\% of the reported value the network alone places as many frames as the causal
GP, not more (the registered statement is not met; Table~\ref{tab:main}b), and its median
absolute error is not below that of hold. The filter does what the network alone does
not: KFX places more frames than the causal GP within 10\% and within 20\% of the value
(\Supp{sec:relerr}), and its median absolute error is below that of hold
(\Supp{sec:filter}). So does the
two-stage estimator, within 20\% and for the median by registered statements
(Table~\ref{tab:main}c), within 10\% as a reported readout (\SuppTab{tab:twopairs}).

% =========================================================================
\subsection{Which diagnostic carries the increment}
\label{sec:res:diag}

BES is light of the heating beam, and the CES value needs the beam. Registered readouts
ask which diagnostic carries the increment and whether it needs a change of the beam
(Table~\ref{tab:carry}; all on the first stage): the pooled network given channels of the
row before (not retrained), networks retrained on the folds of three re-partitions with
changed inputs,
networks that read one component of the BES array, and the filter with one BES summary at
a time.

\begin{table}[tbp]
  \centering
  \caption{What carries the increment (skill with 95\% interval; bold: the interval
  excludes zero). First block: the trained network (first stage) of the pooled
  cross-validation with some channels replaced at test time by the values of the row
  before, or held at their training mean. Second block: networks trained on every fold of
  re-partitions r2--r4 with all fast channels one row old (L), with BES or ECEI only beside
  the CES history (B, E), or without the Mirnov channels (N); the earlier ablations on four
  file-level splits are in \SuppTab{tab:ablation}. Third block: networks on four
  file-level splits, one initialisation each (range of the four point estimates; number of
  splits whose interval is above zero; \SuppTab{tab:components}). Last block: the
  five-parameter filter with a changed drive.}
  \label{tab:carry}
  \footnotesize
  \TabCarry
\end{table}

The information is in the BES value of the present frame: a network trained with all
fast channels one row old has no increment and the first-stage network gains over it, a
network that reads BES and the CES history has the increment and the first-stage network
is not separated from it, and the trained network loses its increment when only its BES
channels are one frame old. A network that reads ECEI and the CES history has no
increment. The Mirnov inputs are not needed: a network without them has the increment and
is not separated from the first-stage network. A trained network of this kind
needs the common mode of the BES array, in which the beam and
the density are not separated, and the filter gets most of its increment from the common
mode too. With the radial difference alone, which attenuates a beam interruption about
eightfold (\Supp{sec:rows}), the filter keeps two thirds of its increment, and the
increment does not need the rows after a blip. Trained for a fixed 40 epochs, the
network keeps its increment, so it does not rest on early stopping. The increment is
therefore carried by the BES value of the present frame; a change of the beam alone is not
needed to explain it, and it is not excluded, because the files carry no beam power.

% =========================================================================
\subsection{Where the gain is}
\label{sec:res:where}

The stored row errors of the network of the pooled cross-validation and of the filter were
read by temperature level, by the rate of edge bursts, by row class and over the error
distribution (registered; Fig.~\ref{fig:where}, Table~\ref{tab:where}; definitions in
\Supp{sec:rows}).

\begin{figure}[tbp]
  \centering
  \includegraphics[width=\linewidth]{fig_where}
  \caption{Where the gain is (numbers in Table~\ref{tab:where}). (a) Skill by thirds of
  the discharges ordered by their median \Ti{}: the network (first stage) over the causal
  GP and over its twin (r2--r9), and KFX over KF0 (r3--r9). (b) Increment of the network
  in rows not flagged by plain value rules, in rows flagged by them (values
  below 150\,eV or above 2.5\,keV, a jump of more than 1\,keV, or low beam emission in the
  30\,ms before) and in the three rows after an absent frame. Filled markers: the interval
  excludes zero.}
  \label{fig:where}
\end{figure}

\begin{table}[tbp]
  \centering
  \caption{Where the gain is (pooled, r2--r9 unless marked; 95\% intervals; bold: the
  interval excludes zero for a skill or a difference, one for a ratio, one half for a share
  of discharges). The increment is the skill of the network (first stage) over its twin.
  Quiet rows have no absent frame, low beam emission or large BES departure behind them;
  rows weighted to resemble rows without a value are weighted by the odds of an unusable
  value predicted from the flags of the five rows before and the fast channels.
  Edge-burst rate: onsets per second of upward departures of the BES channels, none,
  up to 0.719 per second and above; 0.719 is the median of the 360 non-zero rates, and the
  24 discharges whose rate equals it exactly fall in the lower class, which therefore holds
  201 discharges and the upper 159; 28 discharges too short for a rate are in no class. Rows between
  bursts are at least 50\,ms from every onset. Four-split values and the other rows: \SuppTab{tab:pooledrows}.}
  \label{tab:where}
  \footnotesize
  \setlength{\tabcolsep}{2pt}
  \TabWhere
\end{table}

By temperature level, the increment has an interval above zero in the middle and hottest
thirds of the discharges and not in the coldest, for the network and for the filter alike,
while the margin over the causal GP has one in every third (Fig.~\ref{fig:where}a). No
interval was computed for the difference between thirds on the pooled rows, and on single
splits those differences include zero (\Supp{sec:rows}), so a dependence on the level is
not shown. In the coldest third the network places fewer frames within 20\%
of the reported value than the causal GP; KFX does not. After the second stage, on three
re-partitions, the increment and the gain over the network have intervals above zero in
every third, with the smallest point values in the coldest (\SuppTab{tab:twothirds}; reported, not a
registered statement). Classed by the rate of edge
bursts in the BES channels, the increment is established in the rows between bursts and is
not tied to one class: the class with no bursts and the class with frequent bursts are
not separated from each other (Table~\ref{tab:where}); the frequent-burst class is the coldest, and within thirds of the
temperature the classes show no effect of their own (\SuppTab{tab:classthird}). The
increment is in the rows that
plain value rules do not flag, and these rows carry the reduction against the twin; the
margin over the causal GP leans on flagged rows as well (Fig.~\ref{fig:where}b). The
network lowers the 90th percentile of the absolute error but not its median, and has the
lower error in about two thirds of the discharges. The 64 discharges (a tenth) with the
largest gain carry much of it: without them the margin over the causal GP remains, the increment is not
separated from zero (Table~\ref{tab:where}).

\emph{Beam blips.} After a blip of one absent frame the reported \Ti{}, the ECEI level and
the BES level all fall in the first row; over the next rows the deficit of \Ti{} follows
that of the ECEI level, which uses neither the beam nor the spectrometer, and not that of
BES, which recovers within two rows (Fig.~\ref{fig:blips}a, Table~\ref{tab:blips}). On the
raw streams the ECEI level begins to fall while the beam is off and deepens after its
return (Fig.~\ref{fig:blips}b), so a response to the loss of the beam is not excluded in
shape, but the fall is not deeper after two absent frames than after one. In the three
rows after an absent frame the network gains most over its twin (Fig.~\ref{fig:where}b):
the twin sees the longer time step but not the depth of the fall.

\begin{figure}[tbp]
  \centering
  \includegraphics[width=\linewidth]{fig_blips}
  \caption{Beam blips. (a) Event means over blips of one absent frame in the 10\,ms files,
  as ratios to the median of the six rows before (shaded: 95\% interval; dotted: the same
  for control rows without an absent frame). (b) Event medians at 1\,ms resolution through
  the interruptions of the twelve discharges with raw streams, as ratios to the 20\,ms
  before the beam goes off; ECEI over the interruptions kept by a rule fixed beforehand.
  Event sets and definitions: \Supp{sec:rows}.}
  \label{fig:blips}
\end{figure}

\begin{table}[tbp]
  \centering
  \caption{Beam blips (one absent frame unless stated), with 95\% intervals.
  Event sets and definitions: \Supp{sec:rows}.}
  \label{tab:blips}
  \footnotesize
  \begin{tabular}{p{0.68\linewidth}l}
    \toprule
    \multicolumn{2}{l}{\emph{First row after a blip, ratio to the level before (1{,}262 events in 489 discharges)}} \\
    \quad reported \Ti{} & 0.814 $[0.790, 0.835]$ \\
    \quad ECEI level & 0.929 $[0.918, 0.942]$ \\
    \quad BES level & 0.658 $[0.653, 0.663]$ \\
    \multicolumn{2}{l}{\emph{Rows 1 to 8: share of the \Ti{} deficit profile left unexplained}} \\
    \quad by the ECEI deficit & 0.009 $[0.005, 0.032]$ \\
    \quad by the BES deficit & 0.333 $[0.286, 0.376]$ \\
    \quad part of the first-row deficit (0.186) proportional to the BES deficit; one third is 0.062 & 0.030 $[0.001, 0.065]$ \\
    \multicolumn{2}{l}{\emph{Dose}} \\
    \quad first-row deficit of \Ti{}, two absent frames minus one (\DoseTwoEvents{}) & \DoseTi{} \\
    \multicolumn{2}{l}{\emph{Raw streams, 31 interruptions: ECEI level}} \\
    \quad last 3\,ms before the beam returns, ratio (first implementation: 0.969 $[0.945, 1.001]$) & 0.972 $[0.953, 0.999]$ \\
    \quad step at the return & $-0.015$ $[-0.027, +0.012]$ \\
    \quad 10 to 13\,ms after the return minus just after & $-0.018$ $[-0.059, -0.006]$ \\
    \quad first-frame excess of \Ti{} beyond the ECEI shape, share of the level & 0.030 $[0.008, 0.058]$ \\
    \quad the same, share of the first-row deficit & 0.160 $[0.046, 0.292]$ \\
    \quad rotation of the chord, first row, ratio & 0.973 $[0.960, 0.989]$ \\
    \multicolumn{2}{l}{\emph{Network (first stage), rows 1 to 3 after an absent frame}} \\
    \quad skill over its twin & $+0.151$ $[+0.093, +0.215]$ \\
    \quad share within 20\%, minus that of the causal GP & \TTwentyBlip{} \\
    \bottomrule
  \end{tabular}
\end{table}

% =========================================================================
\subsection{Fit or plasma, and conditions of use}
\label{sec:res:use}

Does what the fast channels add describe the plasma or the fit of the frame being
estimated? The next frame is fitted independently, so a correction that describes the
plasma should last into it, one that describes the fit of frame $t$ should not
(registered; read on the first stage). With $d_t$ the network minus its twin, $\beta_h$ is
the regression
coefficient of the reported value of frame $t+h$ minus the twin's estimate for frame $t$
on $d_t$, and the ratio $\rho_h = \beta_h/\beta_0$ measures how much of the correction is
still present $h$ frames later (it can exceed one). Most of the fast-channel correction is
still there in the next frame and
does not fall further over three frames, for the network and the filter
(Fig.~\ref{fig:use}a; \SuppTab{tab:persist}). The ratio $\rho_1$ is a ratio of two
regression coefficients, not the share of the error removed; for the network $\beta_0$ is below one, so its
correction overshoots the reported value of its own frame (for the filter it is near one). A correction from the CES
history alone persists at least fully, so the part of the fast-channel correction that
does not last, about a fifth to a quarter, can be a property of the fit of the frame or a
fluctuation inside it; the assumption that a lasting correction describes the plasma
cannot be tested without fit errors. A bias common to consecutive
frames, through a slowly varying beam or a background shared by neighbouring frames, is
not seen by this test.

\begin{figure}[tbp]
  \centering
  \includegraphics[width=\linewidth]{fig_use}
  \caption{Fit or plasma, and conditions of use. (a) Ratio $\rho_h = \beta_h/\beta_0$, how
  much of a correction is still present in the reported value $h$ frames later, for the
  fast-channel correction of the network and of the filter and for a correction from the
  CES history alone (pooled, r2--r9, for the filter including r2; 95\% intervals; filled:
  the interval excludes one). (b) The networks (first stage) of the
  four splits replayed, not retrained, with the CES value made late by 0 to 4 frames: skill
  of the network over the causal GP (circles, solid) and over its twin (squares,
  dashed), one line per split
  (one initialisation; filled: the interval over the test discharges excludes zero). The
  same with a value visible every fifth frame: \SuppTab{tab:delay}.}
  \label{fig:use}
\end{figure}

The networks of the four splits were replayed with the CES value delayed or thinned
(Fig.~\ref{fig:use}b): with the value one or two frames late or visible every fifth frame,
the margins over the causal GP and over the twin remain on every split. The second stage
was not replayed. On the pooled rows
whose last value is older than 50\,ms nothing is established (Table~\ref{tab:where}).
Real-time BES streaming to a plasma control system is published for
DIII-D~\citep{Kim2026RealTimePlasma}, not for KSTAR. A CES value within about 30\,ms
presupposes on KSTAR a real-time analysis that does not wait for the beam-off frame, the
stated goal of~\citep{Lee2026NeuralNetworkbased}; the CES history of the files is the
offline analysis, so the estimators here are causal in their use of time stamps, not in the
analysis that produced the values they read.

% =========================================================================
\section{Discussion and limitations}
\label{sec:discussion}

The increment is carried by BES, which changes with the beam and with the plasma
(\S\ref{sec:res:diag}). By published slowing-down formulas~\citep{Heidbrink1994BehaviourFast}
a 10\,ms loss of beam heating lowers the temperature by a few per cent at most
(\Supp{sec:rows}), less than the fall of \Ti{} after a blip; a displacement of the plasma
at the blips is reported at KSTAR~\citep{Kim2018EnhancedFast} and cannot be tested without
an equilibrium. The ECEI level is not a temperature (at the edge it depends on density as
well), and no mechanism is named; the fall is not used as evidence on ion heat transport.

\emph{Departure from the registered rule.} The rules registered with the main batch asked
for the pooled statements and the statements on four file-level splits together, both for
the causal network to be the estimator of the paper and for its increment to be stated
as established. On single splits of 96 test discharges, with three initialisations and
the spread between them in the interval, the two statements registered for one training
are not met (each interval excludes zero on two of the four splits); the average of three initialisations,
registered after they had failed, has both margins with zero excluded on every split
(\Supp{sec:splits}). We keep the causal network, because the variant standardized over
the whole discharge cannot run in real time, and give its increment with the pooled
interval; both depart from those rules. The second stage was added under the rule of its
own registration, which made the two-stage estimator the estimator of the paper if its
skills over the network and over the causal GP and its median statement held
(\Supp{sec:twostage}).

\label{sec:limits}
\begin{itemize}\setlength{\itemsep}{1pt}\setlength{\parskip}{0pt}
\item \emph{Target.} The reported C\,VI temperature, not the main-ion temperature; no
per-frame fit error; of a bias that varies within a discharge only a step at beam-off
frames is bounded. The floor rests on an assumption that cannot be tested without fit
errors.
\item \emph{Population.} With values above 3\,keV kept, the increment at the natural
cadence is small for the earlier, whole-discharge variant (\Supp{sec:fast}); the
statement is one about the population without those values.
\item \emph{Beam.} The files record the blips but not the beam power, which would decide
what part of the common mode of BES is the beam.
\item \emph{Level and discharges.} The increment of the network is not established in the
coldest third of the discharges, nor is that of the filter, nor that of the network without
the 64 discharges (a tenth) with the largest reduction or on rows weighted to resemble rows
without a value. Of these readouts only the thirds were repeated after the second stage.
\item \emph{Second stage.} Three re-partitions and one initialisation per network; its
settings come from an earlier batch on single splits whose discharges overlap these; it
reads no information that the network does not read; persistence, delay and the
diagnostic readouts were not repeated on it.
\item \emph{Selection.} The pooled folds re-partition the discharges on which the
recurrent design was adopted; the filter's specification was chosen on one re-partition;
the value rules and the row classes were defined with an earlier training of the
network; ablation, component and replay networks have one initialisation; no untouched
hold-out exists.
\item \emph{Scope and use.} One chord, one campaign, one device; nothing was run on a
plant control system.
\end{itemize}

% =========================================================================
\section{Conclusion}
\label{sec:conclusion}

On the KSTAR edge CES chord selected as the pedestal top, the beam-emission value of the
present frame carries information about the reported ion temperature of that frame beyond
what the CES history gives. A recurrent network against its twin and a five-parameter
filter with and without a BES drive agree on this, and the network given one-frame-old BES
places it in the present frame. For the network, the information is small against an error
that, between two scored frames, is mostly scatter of the reported value; most of it lasts
into the next, independently fitted frame; it is established in the middle and hottest
thirds of the discharges by temperature, not in the coldest. With a second stage fitted on the network's errors, on three
re-partitions, the fast-channel increment keeps an interval above zero, the ratio of the
median error to that of hold has an interval below one, and the skill over the offline GP,
which waits for the next value, has an interval above zero where that of the network alone
is below zero. Use needs the BES value of the present frame
and a CES value at most two frames late; real-time control is a prospect, not a
demonstrated use. The measurements that would change these statements are the beam power,
per-frame fit uncertainties, an equilibrium mapping for the chord, untouched discharges and
further chords.

\section*{Data and code availability}
The MAST data are public (FAIR-MAST, CC BY-SA 4.0); the shot lists, the registration
documents, the analysis code (among it every estimator and reference of this paper)
and all evaluation artifacts are at \url{https://github.com/iseungsang01/ml-intern-revision}. The KSTAR data were made
available to the author through the author's laboratory for this study and cannot be
redistributed by the author; requests for them should be addressed to the Korea
Institute of Fusion Energy.

\section*{Acknowledgments}
The KSTAR data used in this work were provided by the KSTAR team. This work received
no specific funding.

\bibliographystyle{unsrtnat}
\bibliography{refs}

\end{document}
